% Generated draft: requires mathematical and source review.
\documentclass[11pt]{article}
\usepackage[margin=25mm]{geometry}
\usepackage{fontspec}
\setmainfont{Latin Modern Roman}
\usepackage{amsmath,amssymb,mathtools,unicode-math}
\setmathfont{Latin Modern Math}
\usepackage{enumitem,needspace,xurl,hyperref}
\setlength{\parindent}{0pt}
\setlength{\parskip}{5pt}
\setlength{\emergencystretch}{3em}
\newcommand{\N}{\mathbb{N}}
\newcommand{\Z}{\mathbb{Z}}
\newcommand{\Q}{\mathbb{Q}}
\newcommand{\R}{\mathbb{R}}
\newcommand{\C}{\mathbb{C}}
\newcommand{\F}{\mathbb{F}}
\newcommand{\A}{\mathbb{A}}
\newcommand{\PP}{\mathbb{P}}
\newcommand{\E}{\mathbb{E}}
\newcommand{\eps}{\varepsilon}
\newcommand{\dd}{\,\mathrm{d}}
\begin{document}
\section*{UnsolvedMath ID: 30001571}
\textbf{Source problem code:} catalogue entry
\subsection{Definitions and mathematical statement}
Let \(\mathcal C\) denote the source's class of compact complex spaces of K\"ahler type, let \(\mathcal M_{\mathcal C}\) denote the intended many-sorted structure whose sorts are these spaces, and put \(T_{\mathcal C}=\operatorname{Th}(\mathcal M_{\mathcal C})\). The supplied record does not define ``K\"ahler type'' or specify the language and interpretations of this structure. Thus these symbols refer to the intended source construction; the record alone does not determine a complete formal specification of \(T_{\mathcal C}\).

For a complete type \(p\), write \(U(p)\) for its Lascar \(U\)-rank and \(\operatorname{MR}(p)\) for its Morley rank. Write \(p\not\perp q\) to mean that \(p\) and \(q\) are nonorthogonal in the model-theoretic sense. For stationary types in a stable theory, this means that suitable nonforking extensions to a common parameter set have realizations that are not independent over that set. Stability means that no formula has the order property. Nonmultidimensionality is the standard model-theoretic property appearing in the source; the supplied record gives no further definition or parameter conventions for it.

The source uses \emph{almost \(\aleph_1\)-categorical} to mean the conjunction of stability, nonmultidimensionality, and the rank condition below. With that stipulated meaning, the conjecture is
\[
\begin{gathered}
T_{\mathcal C}\text{ is stable and nonmultidimensional},\\
\forall p\;\bigl(U(p)=1\ \Longrightarrow\ \exists q\,
[\operatorname{MR}(q)=1\ \land\ p\not\perp q]\bigr).
\end{gathered}
\]
Here \(p\) and \(q\) range over types of \(T_{\mathcal C}\) under the source's intended conventions. In particular, the witness \(q\) may depend on \(p\); no single common witness is asserted. The record does not specify the allowed parameter sets, tuple arities, use of imaginary sorts, or stationarity conventions. These omissions, together with the unspecified class and language, require clarification before this becomes a fully specified formal statement. No ordinary categoricity assertion is substituted for the source's explicitly stated meaning of ``almost \(\aleph_1\)-categorical.''
\subsection{Short English statement}
The conjecture attributed to Pillay asks whether the theory associated with compact complex spaces of K\"ahler type is stable and nonmultidimensional, with every type of \(U\)-rank one nonorthogonal to some type of Morley rank one. The precise class, structure, and type conventions need clarification from the source. This is a statement-only draft requiring human mathematical review.
\subsection{Sources}
\begin{itemize}
\item Local registry: \texttt{data/unsolvedmath/problems.json}, supplied record with permanent UnsolvedMath ID \(30001571\), exact source problem code \texttt{catalogue entry}, in the Oberwolfach Reports Open Problems collection.
\item Supplied source citation: \emph{Komplexe Analysis} (2011), Oberwolfach Reports, pp.\ 2283--2333, DOI \texttt{10.4171/owr/2010/39}. Supplied source URL: \url{https://doi.org/10.4171/owr/2010/39}. The background describes Oberwolfach Reports as documenting mathematical workshops and their open-problem discussions.
\item Supplied literature reference: ``Pillay, original OWR report,'' \url{https://doi.org/10.4171/owr/2010/39}. The background and literature assessment attribute the extracted conjecture to Pillay; the structured proposer field is unspecified. The proposed-year field is 2010, while the supplied report citation says 2011; both are retained without reconciliation.
\item Literature assessment dated exactly \texttt{2026-08-21}: the supplied curation labels the problem open and states that the identified literature does not resolve Pillay's extracted conjecture and that the report records it as a conjecture. This is supplied dated curation, not an independently verified literature assessment. The background expressly requests independent verification of time-sensitive status and cited claims.
\item The supplied statement-status label is \texttt{corrected\_verified}. Its accompanying verification text reports that mathematical content was checked against the official Oberwolfach report and that notation, context, and question form were normalized without changing the source-level claim. This is a supplied verification claim; no source visit or independent verification is asserted here.
\end{itemize}
\end{document}
